\section{Related work}
\label{sec:related_work}

\paragraph{Petrophysical prediction from logs.}
Well logs provide indirect but dense measurements of rock and fluid properties,
and core-calibrated interpretation is a standard component of petrophysical
analysis \citep{ellis2007well,asquith2004basic}. Data-driven models are now
widely used to map logs to porosity, permeability, saturation, or shale-related
properties \citep{ahmadi2019comparison,wood2020predicting}. These models are
effective when the training distribution is representative and labels are
sufficient, but their usual prediction form is purely supervised. Sparse
measurements in the target well are commonly used as additional features or for
post-hoc calibration, not as explicit constraints in a forward measurement
model.

\paragraph{Sparse compressed sensing.}
Classical compressed sensing reconstructs signals from incomplete linear
measurements by assuming sparsity or compressibility in a known representation
\citep{donoho2006compressed,candes2006robust,candes2006stable}. This view is
appealing for sparse core assimilation because the measurement operator is
known: a row of \(M\) selects the depths where the target was observed. Its
limitation is the structural prior. A fixed basis can be too restrictive for
cross-well petrophysical windows, where changes in depositional setting,
diagenesis, and log-response calibration may create variability that is not
sparse in a universal transform.

\paragraph{Generative priors for inverse problems.}
CSGM replaces the sparse basis with the range of a trained generator
\citep{bora2017compressed}. The unknown signal is represented as \(G(z)\), and
reconstruction searches over \(z\) to minimize measurement mismatch.
\citet{bora2017compressed} established recovery guarantees under set-restricted
eigenvalue conditions and showed that generative priors can require fewer
measurements than sparsity priors in image experiments. Subsequent work studied
the power and limitations of generative compressed sensing
\citep{kamath2020power}, invertible generative priors for inverse problems
\citep{asim2020invertible}, and untrained architectural priors
\citep{ulyanov2018deep}. These works support the central idea that an inverse
problem can be solved by searching a structured generative representation rather
than a fixed sparse basis.

% \paragraph{Conditional generative priors and learned inverse problems.}
% Conditional generative models use auxiliary information to guide generation,
% for example through conditional adversarial models \citep{mirza2014conditional}
% or conditional variational models \citep{sohn2015learning}. In parallel, learned
% inverse-problem methods use neural models as reconstruction maps, regularizers,
% or priors \citep{ongie2020deep,lucas2018using}. CLP-CSGM is related to this
% family because it uses logs as side information, but its inference mechanism is
% different from direct conditional generation: the logs define a latent prior
% \(z_0=h(u)\), while sparse target-well observations remain in the data-fidelity
% term. Thus the conditional information guides the latent search without
% replacing the measurement-consistency objective.

\paragraph{Conditional inverse reconstruction.}
A useful distinction in the inverse-problem literature is between \emph{direct
conditional prediction} and \emph{conditional inverse reconstruction}. In direct
conditional prediction, auxiliary variables are used to predict the output
directly, as in supervised regression, conditional generative modeling, or
latent-code prediction followed by decoding
\citep{mirza2014conditional,sohn2015learning}. In conditional inverse
reconstruction, by contrast, side information does not eliminate the inverse
problem itself; rather, it regularizes the solution while measurement
consistency remains explicit in the objective
\citep{ongie2020deep,lucas2018using}. This distinction is central in our
setting. Dense logs provide contextual information about the target interval,
but sparse core-derived measurements remain the only direct observations of the
property to be reconstructed. A method that simply predicts \(y\) from \(u\), or
even from \([u,b]\), does not preserve this separation of roles.

From this perspective, CLP-CSGM is neither a standard conditional generator nor
a direct latent regressor. Standard conditional generative models use
contextual information to generate samples directly from the conditioning
variable \citep{mirza2014conditional,sohn2015learning}. Such models are
appropriate when the conditional input is intended to determine the output
distribution by itself. Our problem is different: the dense logs \(u\) are
informative but incomplete, while the sparse target-well measurements \(b\)
carry direct evidence about the quantity of interest. For this reason, CLP-CSGM
does not ask the model to synthesize \(y\) from \(u\) alone. Instead, it uses
\(u\) only to define a target-specific latent anchor \(z_0=h(u)\), while the
final estimate is obtained through a measurement-consistent latent inverse
problem.

This formulation is also distinct from standard CSGM
\citep{bora2017compressed}. Classical CSGM searches the generator range by
minimizing only the measurement mismatch. In CLP-CSGM, the latent search is
additionally regularized toward a context-dependent prior predicted from dense
logs. This difference is structural, not merely procedural: the term
\(\lambda\|z-z_0(u)\|_2^2\) remains active throughout the optimization and
changes the reconstructed solution selected from the decoder range. Therefore,
the proposed method should not be interpreted as ``CSGM with a warm start.'' The
logs do not only provide an initialization; they define a persistent latent
regularizer that remains part of the inverse problem during test-time
refinement.

A second relevant distinction is with respect to latent-regression baselines.
One may first learn a latent representation and then regress directly from logs
to latent codes, producing predictions of the form \(G(h(u))\). This already
introduces a learned target manifold, but it still performs \emph{prediction
without test-time assimilation}. Sparse target-well measurements do not alter
the estimate once the latent code has been predicted. In CLP-CSGM, by contrast,
the latent prior and the sparse measurements play complementary roles: the prior
\(z_0=h(u)\) defines a context-dependent center in latent space, while the
measurement term \(\|M G(z)-b\|_2^2\) pulls the estimate toward consistency with
target-well observations. The ablation analysis in this paper confirms that
neither component alone explains the observed gains; the strongest performance
arises from their combination.

The proposed method also differs from fixed-basis sparse recovery. Classical
compressed sensing assumes that the unknown signal is sparse or compressible in
a predetermined transform basis
\citep{donoho2006compressed,candes2006robust,candes2006stable}. That prior can
be too rigid for petrophysical windows affected by cross-well distribution
shift, where geological variability is not naturally aligned with a universal
hand-crafted basis. CLP-CSGM replaces that assumption with a learned generative
range, while preserving the core compressed-sensing principle that observed
entries should appear explicitly in a data-fidelity term.

\paragraph{Position of this work.}
The proposed CLP-CSGM framework is positioned at the intersection of generative
priors, conditional inverse reconstruction, and petrophysical estimation from
sparse core measurements. Its contribution is not a new general theory of
generative compressed sensing, but a task-specific formulation in which dense
logs define a conditional latent prior, sparse target-well measurements define
explicit test-time constraints, and reconstruction is restricted to a learned
range of plausible property windows. This separation of roles distinguishes the
method from direct regression, standard CSGM, and fixed-basis sparse recovery.
